% hw5-function.tex

\newpage
\section{作业 5：函数}

\begin{problem}[函数与等价关系]
  设 $f: X \to Y$ 是满射。
  定义 $X$ 上的二元关系 $R$ 为 $(x, y) \in R$ 当且仅当 $f(x) = f(y)$。
  请证明,
  \begin{enumerate}[(1)]
    \item $R$ 是 $X$ 上的等价关系。
    \item 定义 $h \subseteq (X/R) \times Y$ 为 $h([x]_{R}) = f(x)$。
      请证明, $h$ 是从商集 $X/R$ 到 $Y$ 的函数, 且是满射。
  \end{enumerate}
\end{problem}

\begin{proof}
  \begin{enumerate}[(1)]
    \item 下证 $R$ 是 $X$ 上的等价关系。只需证
      \begin{itemize}
        \item $R$ 是自反的。任取 $x \in X$,
          \[
            f(x) = f(x) \implies (x, x) \in R.
          \]
        \item $R$ 是对称的。任取 $x_{1}, x_{2} \in X$,
          \begin{align*}
            & (x_{1}, x_{2}) \in R \\
            \implies & f(x_{1}) = f(x_{2}) \\
            \implies & f(x_{2}) = f(x_{1}) \\
            \implies & (x_{2}, x_{1}) \in R.
          \end{align*}
        \item $R$ 是传递的。任取 $x_{1}, x_{2}, x_{3} \in X$,
          \begin{align*}
            & (x_{1}, x_{2}) \in R \land (x_{2}, x_{3}) \in R \\
            \implies & f(x_{1}) = f(x_{2}) \land f(x_{2}) = f(x_{3}) \\
            \implies & f(x_{1}) = f(x_{3}) \\
            \implies & (x_{1}, x_{3}) \in R.
          \end{align*}
      \end{itemize}
    \item 下证 $h: X/R \to Y$ 是函数, 且是满射。

		  要证 $h$ 是函数, 只需证
      \begin{itemize}
        \item $\forall S \in X/R\; \exists y \in Y\; h(S) = y$。\\
          设 $X/R = [x]_{R}$。取 $y = f(x)$ 即可。
        \item $\forall S \in X/R\; \exists y_{1}, y_{2} \in Y\;
          ((h(S) = y_{1} \land h(S) = y_{2}) \implies y_{1} = y_{2})$。\\
          根据 $h$ 的定义,
          \[
            \exists x_{1} \in X\; ([x_{1}]_{R} = S \land f(x_{1}) = y_{1}),
          \]
          \[
            \exists x_{2} \in X\; ([x_{2}]_{R} = S \land f(x_{2}) = y_{2}).
          \]
          因此,
          \begin{align*}
            & [x_{1}]_{R} = [x_{2}]_{R} \\
            \implies & (x_{1}, x_{2}) \in R \\
            \implies & f(x_{1}) = f(x_{2}) \\
            \implies & y_{1} = y_{2}
          \end{align*}
      \end{itemize}
      要证 $h$ 是满射, 只需证 $\forall y \in Y\; \exists S \in X/R\; h(S) = y$。\\
      因为 $f$ 是满射, 对于任意 $y \in Y$, 存在 $x_{0}$, 使得 $f(x_{0}) = y$。
      故, 取 $S = [x_{0}]_{R}$ 即可。
  \end{enumerate}
\end{proof}
%%%%%%%%%%%%%%%

\begin{problem}[反函数]
  设 $f: A \to B$ 与 $g: B \to A$ 是函数。
  请证明, \footnote{可以在下周二课程后证明, 当然也可以现在证明。}
  \[
    (f \circ g = I_B \land g \circ f = I_A) \to g = f^{-1}.
  \]
\end{problem}

\begin{proof}
  \begin{intention}
    两个恒等式分别给出 $f$ 的满射性与单射性, 先证明 $f$ 可逆, 再由 $f(g(b)) = b$ 与反函数定义逐点识别 $g(b)$。
  \end{intention}

  设
  \[
    f \circ g = I_B,
    \qquad
    g \circ f = I_A.
  \]

  先证 $f$ 是单射。任取 $a_1, a_2 \in A$, 若
  \[
    f(a_1) = f(a_2),
  \]
  则
  \[
    g(f(a_1)) = g(f(a_2)).
  \]
  由 $g \circ f = I_A$,
  \[
    a_1 = (g \circ f)(a_1) = (g \circ f)(a_2) = a_2.
  \]
  故 $f$ 是单射。

  再证 $f$ 是满射。任取 $b \in B$, 令
  \[
    a = g(b) \in A.
  \]
  则
  \[
    f(a) = f(g(b)) = (f \circ g)(b) = I_B(b) = b.
  \]
  故 $f$ 是满射。

  因而 $f$ 是双射, 故反函数 $f^{-1}: B \to A$ 存在。

  对任意 $b \in B$,
  \[
    f(g(b)) = b.
  \]
  由反函数的定义,
  \[
    f(x) = b \iff x = f^{-1}(b).
  \]
  取 $x = g(b)$, 得
  \[
    g(b) = f^{-1}(b).
  \]
  因而
  \[
    \forall b \in B\; g(b) = f^{-1}(b),
  \]
  即
  \[
    g = f^{-1}.
  \]

  \begin{remark}
    \textbf{重点.} $f \circ g = I_B$ 与 $g \circ f = I_A$ 的作用不同: 前者给满射, 后者给单射。\par
    \textbf{难点.} 证明 $g = f^{-1}$ 之前, 先要说明 $f^{-1}$ 确实存在。\par
    \textbf{易错点.} 不能只用一个恒等式就直接写出 $g = f^{-1}$; 一般必须同时检查左右逆都成立。
  \end{remark}
\end{proof}
